Let \(\mathbb P^{p-1}\) denote real projective space. Define
\[
 \mathcal D(P_X)
 =\bigcup\left\{
   \{[a_k]:1\le k\le m\}:
   \begin{array}{l}
   m\ge p,\ A=[a_1\;\cdots\;a_m]\in\mathbb R^{p\times m},\ \\
   a_k\ne0\text{ and }a_k\notin\operatorname{span}(a_\ell)\text{ for }k\ne\ell,\\
   \xi_1,\ldots,\xi_m\text{ are mutually independent, nondegenerate, and non-Gaussian},\\
   \mathcal L(A\xi)=P_X
   \end{array}
 \right\}.
\]
Thus \(\mathcal D(P_X)\) is defined set-theoretically from the observational law by unioning the unordered projective column sets of all irreducible, all-non-Gaussian independent-component representations. Its asserted uniqueness and equality to the directions of a particular representation are theorem conclusions, not parts of this definition, and the definition does not prescribe an evaluation algorithm.